\documentclass[11pt]{article}

    \usepackage[breakable]{tcolorbox}
    \usepackage{parskip} % Stop auto-indenting (to mimic markdown behaviour)
    

    % Basic figure setup, for now with no caption control since it's done
    % automatically by Pandoc (which extracts ![](path) syntax from Markdown).
    \usepackage{graphicx}
    % Maintain compatibility with old templates. Remove in nbconvert 6.0
    \let\Oldincludegraphics\includegraphics
    % Ensure that by default, figures have no caption (until we provide a
    % proper Figure object with a Caption API and a way to capture that
    % in the conversion process - todo).
    \usepackage{caption}
    \DeclareCaptionFormat{nocaption}{}
    \captionsetup{format=nocaption,aboveskip=0pt,belowskip=0pt}

    \usepackage{float}
    \floatplacement{figure}{H} % forces figures to be placed at the correct location
    \usepackage{xcolor} % Allow colors to be defined
    \usepackage{enumerate} % Needed for markdown enumerations to work
    \usepackage{geometry} % Used to adjust the document margins
    \usepackage{amsmath} % Equations
    \usepackage{amssymb} % Equations
    \usepackage{textcomp} % defines textquotesingle
    % Hack from http://tex.stackexchange.com/a/47451/13684:
    \AtBeginDocument{%
        \def\PYZsq{\textquotesingle}% Upright quotes in Pygmentized code
    }
    \usepackage{upquote} % Upright quotes for verbatim code
    \usepackage{eurosym} % defines \euro

    \usepackage{iftex}
    \ifPDFTeX
        \usepackage[T1]{fontenc}
        \IfFileExists{alphabeta.sty}{
              \usepackage{alphabeta}
          }{
              \usepackage[mathletters]{ucs}
              \usepackage[utf8x]{inputenc}
          }
    \else
        \usepackage{fontspec}
        \usepackage{unicode-math}
    \fi

    \usepackage{fancyvrb} % verbatim replacement that allows latex
    \usepackage{grffile} % extends the file name processing of package graphics
                         % to support a larger range
    \makeatletter % fix for old versions of grffile with XeLaTeX
    \@ifpackagelater{grffile}{2019/11/01}
    {
      % Do nothing on new versions
    }
    {
      \def\Gread@@xetex#1{%
        \IfFileExists{"\Gin@base".bb}%
        {\Gread@eps{\Gin@base.bb}}%
        {\Gread@@xetex@aux#1}%
      }
    }
    \makeatother
    \usepackage[Export]{adjustbox} % Used to constrain images to a maximum size
    \adjustboxset{max size={0.9\linewidth}{0.9\paperheight}}

    % The hyperref package gives us a pdf with properly built
    % internal navigation ('pdf bookmarks' for the table of contents,
    % internal cross-reference links, web links for URLs, etc.)
    \usepackage{hyperref}
    % The default LaTeX title has an obnoxious amount of whitespace. By default,
    % titling removes some of it. It also provides customization options.
    \usepackage{titling}
    \usepackage{longtable} % longtable support required by pandoc >1.10
    \usepackage{booktabs}  % table support for pandoc > 1.12.2
    \usepackage{array}     % table support for pandoc >= 2.11.3
    \usepackage{calc}      % table minipage width calculation for pandoc >= 2.11.1
    \usepackage[inline]{enumitem} % IRkernel/repr support (it uses the enumerate* environment)
    \usepackage[normalem]{ulem} % ulem is needed to support strikethroughs (\sout)
                                % normalem makes italics be italics, not underlines
    \usepackage{soul}      % strikethrough (\st) support for pandoc >= 3.0.0
    \usepackage{mathrsfs}
    

    
    % Colors for the hyperref package
    \definecolor{urlcolor}{rgb}{0,.145,.698}
    \definecolor{linkcolor}{rgb}{.71,0.21,0.01}
    \definecolor{citecolor}{rgb}{.12,.54,.11}

    % ANSI colors
    \definecolor{ansi-black}{HTML}{3E424D}
    \definecolor{ansi-black-intense}{HTML}{282C36}
    \definecolor{ansi-red}{HTML}{E75C58}
    \definecolor{ansi-red-intense}{HTML}{B22B31}
    \definecolor{ansi-green}{HTML}{00A250}
    \definecolor{ansi-green-intense}{HTML}{007427}
    \definecolor{ansi-yellow}{HTML}{DDB62B}
    \definecolor{ansi-yellow-intense}{HTML}{B27D12}
    \definecolor{ansi-blue}{HTML}{208FFB}
    \definecolor{ansi-blue-intense}{HTML}{0065CA}
    \definecolor{ansi-magenta}{HTML}{D160C4}
    \definecolor{ansi-magenta-intense}{HTML}{A03196}
    \definecolor{ansi-cyan}{HTML}{60C6C8}
    \definecolor{ansi-cyan-intense}{HTML}{258F8F}
    \definecolor{ansi-white}{HTML}{C5C1B4}
    \definecolor{ansi-white-intense}{HTML}{A1A6B2}
    \definecolor{ansi-default-inverse-fg}{HTML}{FFFFFF}
    \definecolor{ansi-default-inverse-bg}{HTML}{000000}

    % common color for the border for error outputs.
    \definecolor{outerrorbackground}{HTML}{FFDFDF}

    % commands and environments needed by pandoc snippets
    % extracted from the output of `pandoc -s`
    \providecommand{\tightlist}{%
      \setlength{\itemsep}{0pt}\setlength{\parskip}{0pt}}
    \DefineVerbatimEnvironment{Highlighting}{Verbatim}{commandchars=\\\{\}}
    % Add ',fontsize=\small' for more characters per line
    \newenvironment{Shaded}{}{}
    \newcommand{\KeywordTok}[1]{\textcolor[rgb]{0.00,0.44,0.13}{\textbf{{#1}}}}
    \newcommand{\DataTypeTok}[1]{\textcolor[rgb]{0.56,0.13,0.00}{{#1}}}
    \newcommand{\DecValTok}[1]{\textcolor[rgb]{0.25,0.63,0.44}{{#1}}}
    \newcommand{\BaseNTok}[1]{\textcolor[rgb]{0.25,0.63,0.44}{{#1}}}
    \newcommand{\FloatTok}[1]{\textcolor[rgb]{0.25,0.63,0.44}{{#1}}}
    \newcommand{\CharTok}[1]{\textcolor[rgb]{0.25,0.44,0.63}{{#1}}}
    \newcommand{\StringTok}[1]{\textcolor[rgb]{0.25,0.44,0.63}{{#1}}}
    \newcommand{\CommentTok}[1]{\textcolor[rgb]{0.38,0.63,0.69}{\textit{{#1}}}}
    \newcommand{\OtherTok}[1]{\textcolor[rgb]{0.00,0.44,0.13}{{#1}}}
    \newcommand{\AlertTok}[1]{\textcolor[rgb]{1.00,0.00,0.00}{\textbf{{#1}}}}
    \newcommand{\FunctionTok}[1]{\textcolor[rgb]{0.02,0.16,0.49}{{#1}}}
    \newcommand{\RegionMarkerTok}[1]{{#1}}
    \newcommand{\ErrorTok}[1]{\textcolor[rgb]{1.00,0.00,0.00}{\textbf{{#1}}}}
    \newcommand{\NormalTok}[1]{{#1}}

    % Additional commands for more recent versions of Pandoc
    \newcommand{\ConstantTok}[1]{\textcolor[rgb]{0.53,0.00,0.00}{{#1}}}
    \newcommand{\SpecialCharTok}[1]{\textcolor[rgb]{0.25,0.44,0.63}{{#1}}}
    \newcommand{\VerbatimStringTok}[1]{\textcolor[rgb]{0.25,0.44,0.63}{{#1}}}
    \newcommand{\SpecialStringTok}[1]{\textcolor[rgb]{0.73,0.40,0.53}{{#1}}}
    \newcommand{\ImportTok}[1]{{#1}}
    \newcommand{\DocumentationTok}[1]{\textcolor[rgb]{0.73,0.13,0.13}{\textit{{#1}}}}
    \newcommand{\AnnotationTok}[1]{\textcolor[rgb]{0.38,0.63,0.69}{\textbf{\textit{{#1}}}}}
    \newcommand{\CommentVarTok}[1]{\textcolor[rgb]{0.38,0.63,0.69}{\textbf{\textit{{#1}}}}}
    \newcommand{\VariableTok}[1]{\textcolor[rgb]{0.10,0.09,0.49}{{#1}}}
    \newcommand{\ControlFlowTok}[1]{\textcolor[rgb]{0.00,0.44,0.13}{\textbf{{#1}}}}
    \newcommand{\OperatorTok}[1]{\textcolor[rgb]{0.40,0.40,0.40}{{#1}}}
    \newcommand{\BuiltInTok}[1]{{#1}}
    \newcommand{\ExtensionTok}[1]{{#1}}
    \newcommand{\PreprocessorTok}[1]{\textcolor[rgb]{0.74,0.48,0.00}{{#1}}}
    \newcommand{\AttributeTok}[1]{\textcolor[rgb]{0.49,0.56,0.16}{{#1}}}
    \newcommand{\InformationTok}[1]{\textcolor[rgb]{0.38,0.63,0.69}{\textbf{\textit{{#1}}}}}
    \newcommand{\WarningTok}[1]{\textcolor[rgb]{0.38,0.63,0.69}{\textbf{\textit{{#1}}}}}


    % Define a nice break command that doesn't care if a line doesn't already
    % exist.
    \def\br{\hspace*{\fill} \\* }
    % Math Jax compatibility definitions
    \def\gt{>}
    \def\lt{<}
    \let\Oldtex\TeX
    \let\Oldlatex\LaTeX
    \renewcommand{\TeX}{\textrm{\Oldtex}}
    \renewcommand{\LaTeX}{\textrm{\Oldlatex}}
    % Document parameters
    % Document title
    \title{3904\_lab4}
    
    
    
    
    
    
    
% Pygments definitions
\makeatletter
\def\PY@reset{\let\PY@it=\relax \let\PY@bf=\relax%
    \let\PY@ul=\relax \let\PY@tc=\relax%
    \let\PY@bc=\relax \let\PY@ff=\relax}
\def\PY@tok#1{\csname PY@tok@#1\endcsname}
\def\PY@toks#1+{\ifx\relax#1\empty\else%
    \PY@tok{#1}\expandafter\PY@toks\fi}
\def\PY@do#1{\PY@bc{\PY@tc{\PY@ul{%
    \PY@it{\PY@bf{\PY@ff{#1}}}}}}}
\def\PY#1#2{\PY@reset\PY@toks#1+\relax+\PY@do{#2}}

\@namedef{PY@tok@w}{\def\PY@tc##1{\textcolor[rgb]{0.73,0.73,0.73}{##1}}}
\@namedef{PY@tok@c}{\let\PY@it=\textit\def\PY@tc##1{\textcolor[rgb]{0.24,0.48,0.48}{##1}}}
\@namedef{PY@tok@cp}{\def\PY@tc##1{\textcolor[rgb]{0.61,0.40,0.00}{##1}}}
\@namedef{PY@tok@k}{\let\PY@bf=\textbf\def\PY@tc##1{\textcolor[rgb]{0.00,0.50,0.00}{##1}}}
\@namedef{PY@tok@kp}{\def\PY@tc##1{\textcolor[rgb]{0.00,0.50,0.00}{##1}}}
\@namedef{PY@tok@kt}{\def\PY@tc##1{\textcolor[rgb]{0.69,0.00,0.25}{##1}}}
\@namedef{PY@tok@o}{\def\PY@tc##1{\textcolor[rgb]{0.40,0.40,0.40}{##1}}}
\@namedef{PY@tok@ow}{\let\PY@bf=\textbf\def\PY@tc##1{\textcolor[rgb]{0.67,0.13,1.00}{##1}}}
\@namedef{PY@tok@nb}{\def\PY@tc##1{\textcolor[rgb]{0.00,0.50,0.00}{##1}}}
\@namedef{PY@tok@nf}{\def\PY@tc##1{\textcolor[rgb]{0.00,0.00,1.00}{##1}}}
\@namedef{PY@tok@nc}{\let\PY@bf=\textbf\def\PY@tc##1{\textcolor[rgb]{0.00,0.00,1.00}{##1}}}
\@namedef{PY@tok@nn}{\let\PY@bf=\textbf\def\PY@tc##1{\textcolor[rgb]{0.00,0.00,1.00}{##1}}}
\@namedef{PY@tok@ne}{\let\PY@bf=\textbf\def\PY@tc##1{\textcolor[rgb]{0.80,0.25,0.22}{##1}}}
\@namedef{PY@tok@nv}{\def\PY@tc##1{\textcolor[rgb]{0.10,0.09,0.49}{##1}}}
\@namedef{PY@tok@no}{\def\PY@tc##1{\textcolor[rgb]{0.53,0.00,0.00}{##1}}}
\@namedef{PY@tok@nl}{\def\PY@tc##1{\textcolor[rgb]{0.46,0.46,0.00}{##1}}}
\@namedef{PY@tok@ni}{\let\PY@bf=\textbf\def\PY@tc##1{\textcolor[rgb]{0.44,0.44,0.44}{##1}}}
\@namedef{PY@tok@na}{\def\PY@tc##1{\textcolor[rgb]{0.41,0.47,0.13}{##1}}}
\@namedef{PY@tok@nt}{\let\PY@bf=\textbf\def\PY@tc##1{\textcolor[rgb]{0.00,0.50,0.00}{##1}}}
\@namedef{PY@tok@nd}{\def\PY@tc##1{\textcolor[rgb]{0.67,0.13,1.00}{##1}}}
\@namedef{PY@tok@s}{\def\PY@tc##1{\textcolor[rgb]{0.73,0.13,0.13}{##1}}}
\@namedef{PY@tok@sd}{\let\PY@it=\textit\def\PY@tc##1{\textcolor[rgb]{0.73,0.13,0.13}{##1}}}
\@namedef{PY@tok@si}{\let\PY@bf=\textbf\def\PY@tc##1{\textcolor[rgb]{0.64,0.35,0.47}{##1}}}
\@namedef{PY@tok@se}{\let\PY@bf=\textbf\def\PY@tc##1{\textcolor[rgb]{0.67,0.36,0.12}{##1}}}
\@namedef{PY@tok@sr}{\def\PY@tc##1{\textcolor[rgb]{0.64,0.35,0.47}{##1}}}
\@namedef{PY@tok@ss}{\def\PY@tc##1{\textcolor[rgb]{0.10,0.09,0.49}{##1}}}
\@namedef{PY@tok@sx}{\def\PY@tc##1{\textcolor[rgb]{0.00,0.50,0.00}{##1}}}
\@namedef{PY@tok@m}{\def\PY@tc##1{\textcolor[rgb]{0.40,0.40,0.40}{##1}}}
\@namedef{PY@tok@gh}{\let\PY@bf=\textbf\def\PY@tc##1{\textcolor[rgb]{0.00,0.00,0.50}{##1}}}
\@namedef{PY@tok@gu}{\let\PY@bf=\textbf\def\PY@tc##1{\textcolor[rgb]{0.50,0.00,0.50}{##1}}}
\@namedef{PY@tok@gd}{\def\PY@tc##1{\textcolor[rgb]{0.63,0.00,0.00}{##1}}}
\@namedef{PY@tok@gi}{\def\PY@tc##1{\textcolor[rgb]{0.00,0.52,0.00}{##1}}}
\@namedef{PY@tok@gr}{\def\PY@tc##1{\textcolor[rgb]{0.89,0.00,0.00}{##1}}}
\@namedef{PY@tok@ge}{\let\PY@it=\textit}
\@namedef{PY@tok@gs}{\let\PY@bf=\textbf}
\@namedef{PY@tok@gp}{\let\PY@bf=\textbf\def\PY@tc##1{\textcolor[rgb]{0.00,0.00,0.50}{##1}}}
\@namedef{PY@tok@go}{\def\PY@tc##1{\textcolor[rgb]{0.44,0.44,0.44}{##1}}}
\@namedef{PY@tok@gt}{\def\PY@tc##1{\textcolor[rgb]{0.00,0.27,0.87}{##1}}}
\@namedef{PY@tok@err}{\def\PY@bc##1{{\setlength{\fboxsep}{\string -\fboxrule}\fcolorbox[rgb]{1.00,0.00,0.00}{1,1,1}{\strut ##1}}}}
\@namedef{PY@tok@kc}{\let\PY@bf=\textbf\def\PY@tc##1{\textcolor[rgb]{0.00,0.50,0.00}{##1}}}
\@namedef{PY@tok@kd}{\let\PY@bf=\textbf\def\PY@tc##1{\textcolor[rgb]{0.00,0.50,0.00}{##1}}}
\@namedef{PY@tok@kn}{\let\PY@bf=\textbf\def\PY@tc##1{\textcolor[rgb]{0.00,0.50,0.00}{##1}}}
\@namedef{PY@tok@kr}{\let\PY@bf=\textbf\def\PY@tc##1{\textcolor[rgb]{0.00,0.50,0.00}{##1}}}
\@namedef{PY@tok@bp}{\def\PY@tc##1{\textcolor[rgb]{0.00,0.50,0.00}{##1}}}
\@namedef{PY@tok@fm}{\def\PY@tc##1{\textcolor[rgb]{0.00,0.00,1.00}{##1}}}
\@namedef{PY@tok@vc}{\def\PY@tc##1{\textcolor[rgb]{0.10,0.09,0.49}{##1}}}
\@namedef{PY@tok@vg}{\def\PY@tc##1{\textcolor[rgb]{0.10,0.09,0.49}{##1}}}
\@namedef{PY@tok@vi}{\def\PY@tc##1{\textcolor[rgb]{0.10,0.09,0.49}{##1}}}
\@namedef{PY@tok@vm}{\def\PY@tc##1{\textcolor[rgb]{0.10,0.09,0.49}{##1}}}
\@namedef{PY@tok@sa}{\def\PY@tc##1{\textcolor[rgb]{0.73,0.13,0.13}{##1}}}
\@namedef{PY@tok@sb}{\def\PY@tc##1{\textcolor[rgb]{0.73,0.13,0.13}{##1}}}
\@namedef{PY@tok@sc}{\def\PY@tc##1{\textcolor[rgb]{0.73,0.13,0.13}{##1}}}
\@namedef{PY@tok@dl}{\def\PY@tc##1{\textcolor[rgb]{0.73,0.13,0.13}{##1}}}
\@namedef{PY@tok@s2}{\def\PY@tc##1{\textcolor[rgb]{0.73,0.13,0.13}{##1}}}
\@namedef{PY@tok@sh}{\def\PY@tc##1{\textcolor[rgb]{0.73,0.13,0.13}{##1}}}
\@namedef{PY@tok@s1}{\def\PY@tc##1{\textcolor[rgb]{0.73,0.13,0.13}{##1}}}
\@namedef{PY@tok@mb}{\def\PY@tc##1{\textcolor[rgb]{0.40,0.40,0.40}{##1}}}
\@namedef{PY@tok@mf}{\def\PY@tc##1{\textcolor[rgb]{0.40,0.40,0.40}{##1}}}
\@namedef{PY@tok@mh}{\def\PY@tc##1{\textcolor[rgb]{0.40,0.40,0.40}{##1}}}
\@namedef{PY@tok@mi}{\def\PY@tc##1{\textcolor[rgb]{0.40,0.40,0.40}{##1}}}
\@namedef{PY@tok@il}{\def\PY@tc##1{\textcolor[rgb]{0.40,0.40,0.40}{##1}}}
\@namedef{PY@tok@mo}{\def\PY@tc##1{\textcolor[rgb]{0.40,0.40,0.40}{##1}}}
\@namedef{PY@tok@ch}{\let\PY@it=\textit\def\PY@tc##1{\textcolor[rgb]{0.24,0.48,0.48}{##1}}}
\@namedef{PY@tok@cm}{\let\PY@it=\textit\def\PY@tc##1{\textcolor[rgb]{0.24,0.48,0.48}{##1}}}
\@namedef{PY@tok@cpf}{\let\PY@it=\textit\def\PY@tc##1{\textcolor[rgb]{0.24,0.48,0.48}{##1}}}
\@namedef{PY@tok@c1}{\let\PY@it=\textit\def\PY@tc##1{\textcolor[rgb]{0.24,0.48,0.48}{##1}}}
\@namedef{PY@tok@cs}{\let\PY@it=\textit\def\PY@tc##1{\textcolor[rgb]{0.24,0.48,0.48}{##1}}}

\def\PYZbs{\char`\\}
\def\PYZus{\char`\_}
\def\PYZob{\char`\{}
\def\PYZcb{\char`\}}
\def\PYZca{\char`\^}
\def\PYZam{\char`\&}
\def\PYZlt{\char`\<}
\def\PYZgt{\char`\>}
\def\PYZsh{\char`\#}
\def\PYZpc{\char`\%}
\def\PYZdl{\char`\$}
\def\PYZhy{\char`\-}
\def\PYZsq{\char`\'}
\def\PYZdq{\char`\"}
\def\PYZti{\char`\~}
% for compatibility with earlier versions
\def\PYZat{@}
\def\PYZlb{[}
\def\PYZrb{]}
\makeatother


    % For linebreaks inside Verbatim environment from package fancyvrb.
    \makeatletter
        \newbox\Wrappedcontinuationbox
        \newbox\Wrappedvisiblespacebox
        \newcommand*\Wrappedvisiblespace {\textcolor{red}{\textvisiblespace}}
        \newcommand*\Wrappedcontinuationsymbol {\textcolor{red}{\llap{\tiny$\m@th\hookrightarrow$}}}
        \newcommand*\Wrappedcontinuationindent {3ex }
        \newcommand*\Wrappedafterbreak {\kern\Wrappedcontinuationindent\copy\Wrappedcontinuationbox}
        % Take advantage of the already applied Pygments mark-up to insert
        % potential linebreaks for TeX processing.
        %        {, <, #, %, $, ' and ": go to next line.
        %        _, }, ^, &, >, - and ~: stay at end of broken line.
        % Use of \textquotesingle for straight quote.
        \newcommand*\Wrappedbreaksatspecials {%
            \def\PYGZus{\discretionary{\char`\_}{\Wrappedafterbreak}{\char`\_}}%
            \def\PYGZob{\discretionary{}{\Wrappedafterbreak\char`\{}{\char`\{}}%
            \def\PYGZcb{\discretionary{\char`\}}{\Wrappedafterbreak}{\char`\}}}%
            \def\PYGZca{\discretionary{\char`\^}{\Wrappedafterbreak}{\char`\^}}%
            \def\PYGZam{\discretionary{\char`\&}{\Wrappedafterbreak}{\char`\&}}%
            \def\PYGZlt{\discretionary{}{\Wrappedafterbreak\char`\<}{\char`\<}}%
            \def\PYGZgt{\discretionary{\char`\>}{\Wrappedafterbreak}{\char`\>}}%
            \def\PYGZsh{\discretionary{}{\Wrappedafterbreak\char`\#}{\char`\#}}%
            \def\PYGZpc{\discretionary{}{\Wrappedafterbreak\char`\%}{\char`\%}}%
            \def\PYGZdl{\discretionary{}{\Wrappedafterbreak\char`\$}{\char`\$}}%
            \def\PYGZhy{\discretionary{\char`\-}{\Wrappedafterbreak}{\char`\-}}%
            \def\PYGZsq{\discretionary{}{\Wrappedafterbreak\textquotesingle}{\textquotesingle}}%
            \def\PYGZdq{\discretionary{}{\Wrappedafterbreak\char`\"}{\char`\"}}%
            \def\PYGZti{\discretionary{\char`\~}{\Wrappedafterbreak}{\char`\~}}%
        }
        % Some characters . , ; ? ! / are not pygmentized.
        % This macro makes them "active" and they will insert potential linebreaks
        \newcommand*\Wrappedbreaksatpunct {%
            \lccode`\~`\.\lowercase{\def~}{\discretionary{\hbox{\char`\.}}{\Wrappedafterbreak}{\hbox{\char`\.}}}%
            \lccode`\~`\,\lowercase{\def~}{\discretionary{\hbox{\char`\,}}{\Wrappedafterbreak}{\hbox{\char`\,}}}%
            \lccode`\~`\;\lowercase{\def~}{\discretionary{\hbox{\char`\;}}{\Wrappedafterbreak}{\hbox{\char`\;}}}%
            \lccode`\~`\:\lowercase{\def~}{\discretionary{\hbox{\char`\:}}{\Wrappedafterbreak}{\hbox{\char`\:}}}%
            \lccode`\~`\?\lowercase{\def~}{\discretionary{\hbox{\char`\?}}{\Wrappedafterbreak}{\hbox{\char`\?}}}%
            \lccode`\~`\!\lowercase{\def~}{\discretionary{\hbox{\char`\!}}{\Wrappedafterbreak}{\hbox{\char`\!}}}%
            \lccode`\~`\/\lowercase{\def~}{\discretionary{\hbox{\char`\/}}{\Wrappedafterbreak}{\hbox{\char`\/}}}%
            \catcode`\.\active
            \catcode`\,\active
            \catcode`\;\active
            \catcode`\:\active
            \catcode`\?\active
            \catcode`\!\active
            \catcode`\/\active
            \lccode`\~`\~
        }
    \makeatother

    \let\OriginalVerbatim=\Verbatim
    \makeatletter
    \renewcommand{\Verbatim}[1][1]{%
        %\parskip\z@skip
        \sbox\Wrappedcontinuationbox {\Wrappedcontinuationsymbol}%
        \sbox\Wrappedvisiblespacebox {\FV@SetupFont\Wrappedvisiblespace}%
        \def\FancyVerbFormatLine ##1{\hsize\linewidth
            \vtop{\raggedright\hyphenpenalty\z@\exhyphenpenalty\z@
                \doublehyphendemerits\z@\finalhyphendemerits\z@
                \strut ##1\strut}%
        }%
        % If the linebreak is at a space, the latter will be displayed as visible
        % space at end of first line, and a continuation symbol starts next line.
        % Stretch/shrink are however usually zero for typewriter font.
        \def\FV@Space {%
            \nobreak\hskip\z@ plus\fontdimen3\font minus\fontdimen4\font
            \discretionary{\copy\Wrappedvisiblespacebox}{\Wrappedafterbreak}
            {\kern\fontdimen2\font}%
        }%

        % Allow breaks at special characters using \PYG... macros.
        \Wrappedbreaksatspecials
        % Breaks at punctuation characters . , ; ? ! and / need catcode=\active
        \OriginalVerbatim[#1,codes*=\Wrappedbreaksatpunct]%
    }
    \makeatother

    % Exact colors from NB
    \definecolor{incolor}{HTML}{303F9F}
    \definecolor{outcolor}{HTML}{D84315}
    \definecolor{cellborder}{HTML}{CFCFCF}
    \definecolor{cellbackground}{HTML}{F7F7F7}

    % prompt
    \makeatletter
    \newcommand{\boxspacing}{\kern\kvtcb@left@rule\kern\kvtcb@boxsep}
    \makeatother
    \newcommand{\prompt}[4]{
        {\ttfamily\llap{{\color{#2}[#3]:\hspace{3pt}#4}}\vspace{-\baselineskip}}
    }
    

    
    % Prevent overflowing lines due to hard-to-break entities
    \sloppy
    % Setup hyperref package
    \hypersetup{
      breaklinks=true,  % so long urls are correctly broken across lines
      colorlinks=true,
      urlcolor=urlcolor,
      linkcolor=linkcolor,
      citecolor=citecolor,
      }
    % Slightly bigger margins than the latex defaults
    
    \geometry{verbose,tmargin=1in,bmargin=1in,lmargin=1in,rmargin=1in}
    
    

\begin{document}
    
    \maketitle
    
    

    
    \#

PID Temperature Control with an Arduino

    

    

    

    

    

    

    

    

    

    

    

    

    

    \#\#\#

Phy 3902 Laboratory Experiment 4

    

    

    

    

    

    

    

    

    

    

    

    

    

    \#\#\#

Michelle Johnstone and Patrick Geraghty

    \#\#\#

University of Ottawa

    \#\#\#

November 3, 2025

    \hypertarget{abstract}{%
\subsection{Abstract}\label{abstract}}

    As winter approaches, one of the most delicate balancing acts in a
household is the temperature and the hydro bill. Most modern electrical
heating systems operate periodically depending on input from a
thermistor whose resistance depends on the temperature of its
surroundings. When the temperature sensed by the thermistor goes outside
of a set accepted range, the heater will turn on until the desired
temperature is reached. This is most efficient in an isolated system;
however, it is common to adjust the output of the heating element to a
percentage of its maximum to avoid overshooting the desired temperature
and causing any issues with overheating. In cases where the maintained
temperature is important for a closed or open system, a proportional,
integral, and derivative (PID) algorithm is often the most efficient way
to do so. PID temperature control allows for the heating element to
react to changes in the environment and avoid significant impact to the
system. As an example, many commercial fridges use PID temperature
control as a way to make up for the door being constantly opened and
closed as it is more energy efficient than the on-off model common to
household fridges. Through this experiment it was determined that a PID
algorithm with a basic heating element and thermistor setup controlled
by a wallet-friendly Arduino can maintain a temperature with a standard
deviation of \(XX^{\circ }C\) and a percent error of \(XX\%\). These
results were used to calculate the amount of power supplied to the
heater required to maintain different temperatures to estimate the
specific heat of the metal blocks used in the experiment. The specific
heat was determined to be \(\frac{J}{g^{\circ}C}\) for aluminum and
\(\frac{J}{g^{\circ}C}\) for brass.

    \hypertarget{introduction}{%
\subsection{Introduction}\label{introduction}}

    The standard on-off model of temperature control needs to be relatively
inaccurate with a large hysteresis, otherwise the life span of circuitry
decreases proportionally with the temperature range. The ability of the
standard on-off model to maintain a set temperature is severely
diminished because of this and they tend to work better for realistic
thermally isolated systems which don't need a large amount of power to
maintain the set temperature. However, in this experiment the system is
open to it's environment and sudden changes in temperature from external
interaction, this is the benefit of PID control, the system is able to
react to these changes and adjust accordingly. The first step is to make
sure the power provided to the heater is proportional to the amount
needed for the set temperature, this helps to prevent the heater from
overshooting the temperature as the power supplied to the heater
decreases as it approaches the set temperature and increases as it drops
below the desired range. The percentage of power applied is calculated
using the error \(E\) as the difference between the desired and measured
temperatures over the width of the allowed temperature band,
\[E = \frac{T_{\text{set}} - T_{\text{measured}}}{\text{band}}, \tag{1}\]
which was used to control the amount of power supplied to the heater as
a percentage of the maximum amount of power
\[ \text{Power} = ( T_{\text{offset}} + k_p E ) P_{\text{max}}. \tag{2}\]
The factor \(k_p\) was added to control the effects of the error
calculation and control the speed with which the system reacts. While
this modification continues to produce a notable degree of error, it can
be observed that the temperature-time graph follows a decaying
sinusoidal function permits further refinement by integrating over an
interval that is twice the time required for the temperature to increase
by an amount equivalent to the band amplitude. The integral of the error
over this time interval \[I = \int {E } dt\, \tag{3}\] can be used in
conjunction with a proportionality constant \(k_i\) to avoid
overcorrecting the power and introducing new unwanted temperature
oscillations. Adding this factor to equation (2) gives the following
formula for a steadier code for maintaining the set temperature
\[\text{Power} = ( T_{\text{offset}} + k_p E + k_i I) P_{\text{max}}. \tag{4}\]
The final adjustment is to control the temperature when reacting to
external stimuli using derivative control. If the temperature of the
environment or the block suddenly spikes or drops, this factor will
allow the heater to react proportionally to the rate of change of the
temperature of the block. This derivative \(D\) can be found by dividing
the negative of the change in error by the time interval over which this
change is measured, the action time of the code, using the following
equation
\[D = - \frac{dE}{dx} = - \frac{\Delta{E}}{\Delta{x}}. \tag{5}\] This
value can be used in conjunction with an associated proportionality
constant \(k_d\) and introduced to equation (4) to produce the final
equation for maintaining the set temperature
\[\text{Power} = ( T_{\text{offset}} + k_p E + k_i I + k_d D) P_{\text{max}}. \tag{6}\]
While the sum of all three of these terms is very beneficial for a fast
and highly variable system, in this experimental setup it is possible
that the derivative control can cause any noise in the temperature
sensor to be amplified so it may be beneficial to use a low value or
zero for the proportionality constant \(k_d\). The final step was to
determine the specific heat capacity of two metal blocks, this was done
by using the PID code at 3 different set temperatures to determine the
amount of power in watts required to maintain those temperatures and
thus the heat loss from the block to the surroundings. Then by writing a
program to supply 100\% power to the heater until a maximum temperature
was reached delivered a reading for the amount of time \(t\) it took to
raise the temperature of the block by 1°C with 12V of power supplied.
Using the equation for specific heat \(c\),
\[ c = \frac{(P_{\text{heater}} - P_{\text{loss}}) * \Delta{t}}{m * \Delta{T}}, \tag{7}\]
where \(m\) is the mass of the block in grams. For this experiment a
MOSFET determines how much power is supplied to the heater based on a
duty cycle out of 255 which maps the 5V PWM source to the additional 12V
power source connected to the heater meaning the power supplied to the
heater was determined using
\[ P = \frac{V^2_{\text{supply}}}{R_{\text{heater}}} * \frac{L}{255}. \tag{8}\]
where \(L\) is the duty cycle of the MOSFET. The error throughout the
experiment was calculated using the traditionally accepted model
\[ \text{\%error} = \frac{\text{|Experimental - Theoretical|}}{\text{Theoretical}} * 100\%. \tag{9}\]

    \hypertarget{methods}{%
\subsection{Methods}\label{methods}}

    The first step of this experiment is to establish a code to trigger the
heating element to turn on and off with the temperature reading of the
thermistor. The initial code to expand on turned on the heating element
at maximum power until the thermistor registered a set temperature of
65°C at which point it turned the element off until it reached the low
temperature limit. The temperature was recorded roughly every second and
the temperature read from the thermistor was graphed as a function of
the elapsed time, the shape of this graph was expected to be a positive
linear trend upward with a peak generously above the set temperature
with a steeper decline to the low end of the passband indicating the
circuit functioned as expected. The setup for this experiment was the
same throughout as the only change required to turn an on-off
temperature control into a PID control is a change in code.

    Figure 1: Wire diagram of the set up used throughout the entirety of the
experiment.

    After a baseline reactivity of the system was established using the
on-off control code, equation (2) was used to control the output power
to the heater. This allowed for the power to the heater to be modulated
based on the error between the actual block temperature and the set
temperature and for the hysteresis to be better enforced and lessen the
likelihood of the temperature exceeding the band. The proportionality
constant \(k_p\) was set to 0.1 and slowly increased until a noticeable
benefit was noticed. While this improved upon the ability of the heater
to keep the block within the hysteresis, it maintained a quite large gap
from the set temperature and overshot the goal with each oscillation. To
improve upon this the code, the integral control is implemented to
minimize the gap between the measured and set temperatures to increase
the rate of decay of the sinusoidal shape of meausred temperature. The
reaction time of the system is the time it takes for the temperature to
change by the value of the band, the integral action time should be
twice this value to give the system time to stabilize. Using a low
proportionality constant \(k_i\) \textasciitilde0.0001 further helps to
minimize speed of system reactions, the constant can be slowly increased
until the desired effects are observed without overcorrection.

The final code modification for the PID temperature control system is
the derivative control, this is best to ensure a system can combat
sudden changes in environmental temperature. In the lab environment this
can be simulated, however, after adding the derivative factor for a
fairly stable surrounding environment or a slow changing system could
increase the system overreacting and overshooting the desired
temperature due to noise from the thermistor reading and amplify
measurement error which would induce oscillations cancelled out by the
integral factor. The derivative factor is especially useful for systems
with rapidly changing dynamics like motors as it dampens these changes
in the system. In order to minimize unwanted effects, the
proportionality constant \(k_d\) was set to 0.01 and decreased until the
system stopped overreacting to noise in the measurements.

    The temperature-time graphs were analyzed at 2 different set
temperatures to determine the voltage required to maintain each desired
value using the duty cycle of the MOSFET at the stabilized end of the
graph by plugging it into equation (8) to find the heat loss from the
block. Next the specific\_heat\_measurement.ino code which merely
applied full power to the heating element until a max set temperature
was reached from which the values were obtained for the amounts of time
and power required to raise the temperature of the block by 1°C.
Equation (7) used these values to determine the specific heat of both
blocks and the alternate temperatures were used to help estimate the
accuracy and observe how the specific heat changes between low and high
temperatures.

    \hypertarget{results-and-discussion}{%
\subsection{Results and Discussion}\label{results-and-discussion}}

    For the first part of the experiment that involved verifying all
components of the circuit were functioning as expected the temperature
of the simple on-off control code was graphed as a function of time. The
resulting shape of the graph was as expected, a steady linear increase
in temperature with a curved peak slightly above the set temperature of
65°C indicating that this strategy of heating tends to overshoot the set
temperature. After the peak there was a steeper linear decrease in
temperature until it reached the low end of the allowed temperature band
indicating the block lost the heat faster than it gained it which points
to this being a poor method of temperature control for an open system.
This data functioned as a baseline for the required improvements to
maintain a set temperature as well as a reference for the improved
performance and efficiency of a PID filter versus a traditional on-off
temperature control.

    Figure 2: The graph of temperature as a function of time for the simple
on-off circuit with a band of 3.

    After the code was improved upon, the anticipated effects were observed
in the temperature time graphs for each of the modifications. When
equation (2) was used to determine the amount of power supplied to the
heater in terms of the duty cycle of the MOSFET, it was evident that
this method of tempurature control was already significantly more
accurate than the on-off model. The temperature of the block was within
the band significantly longer and it overshot the upper limit of the
band by much less. This is also where the decaying sinusoidal nature of
these style of systems can start to be observed as the proper use of
error in the calculation for power output ensures the system reacts
according to surroundings with significantly less delay.

    Figure 3: The temperature-time graph of the proportionally controlled
heater circuit with a band of 1 in accordance with equation (2).

    Due to the constant offset the temperature was still not quite accurate
to the desired results as the heater maintained a non-zero power input
even when the target temperature was reached resulting in a steady state
solution just higher or lower than the set point. This is where the
integral portion came into play as it further ensured that the use of
the error value in the function has a long-term effect as well as an
immediate one. The variable ensured the total error decreased over each
action time resulting in a very accurate steady state temperature. This
effect is observed in the following graph as the amplitude of the
function decays significantly every 50 seconds which was the value of
the action time.

    Figure 4: The temperature-time graph of the proportionally and
integrally controlled heater circuit with a band of 1 in accordance with
equation (4).

    The system, however, was still very slow to react to external stimuli
whereas a notable observation with the inclusion of the derivative
control was that when there was a sharp drop in temperature observed,
the system reacted significantly more quickly than without. During one
run of the code the thermistor was bumped out of the block resulting in
the program observing a rapid decrease in temperature, this was the only
time during the experiment the heater went to 100\% power output without
it being explicitly set within the code. At the time of this run the
derivative constant was set to 0.001 proving the impact that even
minimized variables can have on a reactive system. On the other hand,
for the majority of the runs where \(k_d\) was non-zero, the
differentiation variable emphasized the noise, or high frequency
changes, likely caused by sensor fluctuations inherent to the thermistor
or digitization of the readings.

    Figure 5: The temperature-time graph of the proportionally, integrally,
and derivatively controlled heater circuit with a band of 1 in
accordance with equation (6).

    While the results and graphs obtained for the temperature control
algorithms all had very little deviation from the expected results, the
same could not be said for the specific heat capacity calculation. For
the aluminum block which had a mass of 24.64 ± 0.005g and took about 4.4
seconds to heat up by 1°C, the specific heat was calculated to be around
\(1.8 \frac{J}{g^{\circ}C}\) whereas the accepted value is
\(0.9 \frac{J}{g^{\circ}C}\). It was announced beforehand that the
results for aluminum tend to come out higher than the expected result,
however, this 100\% error indicates either that the calculation or the
measurements were incredibly far off. The measurements for the brass
block yield similar issues, it had a mass of 72.35 ± 0.005g and took
about 5.3 seconds to heat up by 1°C, the specific heat was calculated to
be around \(0.8 \frac{J}{g^{\circ}C}\) whereas the accepted value is
\(0.4 \frac{J}{g^{\circ}C}\).

    Figure 6: ********A graph related to specific heat capacity???********

    Perhaps getting a wider range of values would have helped in determining
where in the setup the error was, however, due to complications with
broken hardware throughout the experiment, the expected timeline for the
experiment was significantly altered. The most likely culprit seemed to
be an error in the code for the scaling, recording, or calculation of
the power value at either the PWM or MOSFET level or in the code
determining joule output in the specific heat measurement calculation or
amplification of noise before the calculation.

    \hypertarget{conclusion}{%
\subsection{Conclusion}\label{conclusion}}

    In conclusion, the Arduino-based PID temperature control system
effectively demonstrated progressive improvements in stability and
response across the ON--OFF, P, PI, and PID configurations. The PI
controller achieved the most stable performance, maintaining the target
temperature with the smallest standard deviation and least noise
reactivity, validating its theoretical advantages. However, the
experimentally determined specific heat capacities for aluminum and
brass were significantly higher than accepted values, primarily due to
systematic errors. These could have included code-related issues such as
the chosen constants or improper power scaling of the PWM signals,
circuitry losses through the MOSFET and wiring, limited thermal
insulation, and sensor inaccuracies from the thermistor. Had there been
more time for the experiment it could have proven beneficial to
incorporate direct voltage and current measurements for accurate power
determination, improved sensor calibration, and refined control code to
eliminate offset bias. Despite these discrepancies, the experiment
successfully demonstrated the principles and effectiveness of PID-based
temperature regulation.

    \hypertarget{references}{%
\subsection{References}\label{references}}

    Webpages:

Michel Godin;
https://uottawa.brightspace.com/content/enforced/523665-2259FS0268421B00/PHY\%203902\_LED\_E\_2023.pdf
(Accessed on October 20, 2025)

    \begin{tcolorbox}[breakable, size=fbox, boxrule=1pt, pad at break*=1mm,colback=cellbackground, colframe=cellborder]
\prompt{In}{incolor}{ }{\boxspacing}
\begin{Verbatim}[commandchars=\\\{\}]

\end{Verbatim}
\end{tcolorbox}


    % Add a bibliography block to the postdoc
    
    
    
\end{document}
